% created by make_collection.py
% manual adjustments:
% - add space to abbreviaton vla d'amore
% - set abbrwidth 6em
% - set editor, ORCID, ARK
% - from front matter of HerEy 115
%   - add abbreviation for fl
%   - add description
%   - include pdf, add option openright
\documentclass[tocdir=../../tmp/B4,abbrwidth=6em]{ees}

\begin{document}

\title{[collection – add custom title page]}
\date{2026-02-28}
\license{cc-by-sa-4.0}
\def\MetadataLilypondVersion{2.24.4}
\def\MetadataEESToolsVersion{v2026.01.0}
\def\MetadataEditor{Wolfgang Esser-Skala}
\def\MetadataOrcid{0000-0002-7350-4045}
\def\MetadataArk{68748/e1jleb4}
\repository{edition-esser-skala/eybler-sacred-music}
\version{v2026.02.0}
\checksum{98058be7b6cd5ed5265d62eeed13efb4dae8333d}
\def\MetadataQRCode{\relax}
\eesTitlePage

\chapter{Critical Report}

In general, this edition closely follows the respective principal source.
Any changes that were introduced by the editor are indicated
by italic type (lyrics, dynamics and directives), parentheses
(expressive marks and bass figures) or dashes (slurs and ties).
Accidentals are used according to modern conventions.
For further details, consult the Editorial Guidelines
available on the Edition’s webpage.

\section{Abbreviations}

\begin{abbreviations}
\abbr{A}{alto}
\abbr{B}{bass}
\abbr{S}{soprano}
\abbr{T}{tenor}
\abbr{b}{basses}
\abbr{clno}{clarion}
\abbr{cl}{clarinet}
\abbr{cor}{horn}
\abbr{fag}{bassoon}
\abbr{fl}{flute}
\abbr{ob}{oboe}
\abbr{org}{organ}
\abbr{timp}{timpani}
\abbr{trb}{trombone}
\abbr{vla d’amore}{Viola d’amore}
\abbr{vla}{viola}
\abbr{vlc}{violoncello}
\abbr{vl}{violin}
\end{abbreviations}

\section{Te summe Jesu · HerEy 49}

\begin{xltabular}{\linewidth}{@{} >\itshape l X}
genre & gradual \\
festival & de SS. Nomine Jesu et de Tempore \\
scoring & S, A, T, B (coro), 2 ob, 2 fag, 2 clno (D), timp (D–A), 2 vl, vla, vlc, b, org \\
\end{xltabular}

\begin{sources}

\sourceitem%
  {A1}%
  {A-Ws}%
  {569 (2)}%
  {autograph manuscript (principal source)}%
  {1809}%
  {}%
  {public domain}%
  {}%
  {full score; 16 pages}

\end{sources}

\begin{xltabular}{\linewidth}{ll X}
\toprule
\itshape Bar & \itshape Staff & \itshape Description \\
\midrule \endhead
– & – & This work is a revised version of the second movement
        of \textit{Hæc est dies} (HerEy~89). See the critical report
        of the latter work for details. \\

\bottomrule
\end{xltabular}

\textlt{\textit{Lyrics}\\Te summe Jesu fontem amoris\\ beati coelites adorent,\\ et triumphali cantico laudent\\ per cuncta saecula.}

\section{Iste est, qui ante Deum · HerEy 69}

\begin{xltabular}{\linewidth}{@{} >\itshape l X}
genre & gradual \\
festival & de Sancti Josephi \\
scoring & S, A, T, B (coro), 2 ob, [2 cl (\flat B)], 2 fag, 2 clno (\flat B), timp (\flat B–F), 2 vl, vla, vlc, b, org \\
\end{xltabular}

\begin{sources}

\sourceitem%
  {A1}%
  {A-Ws}%
  {Cod. 571/3}%
  {autograph manuscript (principal source)}%
  {1807}%
  {}%
  {public domain}%
  {}%
  {full score (8 pages); cl 1/2 on one extra page}

\end{sources}

\begin{xltabular}{\linewidth}{ll X}
\toprule
\itshape Bar & \itshape Staff & \itshape Description \\
\midrule \endhead
–  & cl & Clarinets play unison with oboes and thus should be
          regarded as alternative instruments. \\
13 & T  & 4th \quarterNote\ in \A1: d′4 \\

\bottomrule
\end{xltabular}

\textlt{\textit{Lyrics}\\Iste est qui ante Deum magnas virtutes operatus est,\\ et de omni corde suo laudavit Dominum.\\ (Cantus Index 003426)}

\clearpage
\section{Populum humilem salvum · HerEy 72}

\begin{xltabular}{\linewidth}{@{} >\itshape l X}
genre & gradual \\
festival & de Tempore \\
scoring & S, A, T, B (coro), 2 ob, 2 cl (\flat B), 2 fag, 2 cor (\flat B), 2 vl, vla, vlc, b, org \\
\end{xltabular}

\begin{sources}

\sourceitem%
  {A1}%
  {A-Ws}%
  {Cod. 571/17}%
  {autograph manuscript (principal source)}%
  {1829-09-05}%
  {}%
  {public domain}%
  {}%
  {full score; 8 pages}

\end{sources}

\begin{xltabular}{\linewidth}{ll X}
\toprule
\itshape Bar & \itshape Staff & \itshape Description \\
\midrule \endhead
46 & cl 2 & 3rd \quarterNote\ missing in \A1 \\

\bottomrule
\end{xltabular}

\textlt{\textit{Lyrics}\\Populum humilem salvum facies, Domine,\\ et oculos superborum humiliabis.\\ Quoniam quis Deus praeter te, Domine?\\ (\bibleverse{Ps}(18/17:28,32))}

\section{Alleluia confitemini · HerEy 73}

\begin{xltabular}{\linewidth}{@{} >\itshape l X}
genre & gradual \\
festival & in Vigilia Pentecostes \\
scoring & S, A, T, B (coro), b, org \\
\end{xltabular}

\begin{sources}

\sourceitem%
  {A1}%
  {A-Ws}%
  {Cod. 571/21}%
  {autograph manuscript (principal source)}%
  {}%
  {}%
  {public domain}%
  {}%
  {full score; 2 pages}

\end{sources}


\textlt{\textit{Lyrics}\\Alleluia.\\ Confitemini Domino, quoniam bonus:\\ quoniam in saeculum misericordia eius.\\ (\bibleverse{Ps}(106/105:1))\\[1ex] Laudate Dominum, omnes gentes,\\ collaudate eum, omnes populi.\\ Quoniam confirmata est super nos misericordia eius,\\ et veritas Domini manet in aeternum.\\ (\bibleverse{Ps}(117/116:))}

\section{Nos populus tuus · HerEy 76}

\begin{xltabular}{\linewidth}{@{} >\itshape l X}
genre & offertorium \\
festival & de Tempore \\
scoring & S, A, T, B (coro), 2 ob, 2 fag, 2 clno (C), timp (C–G), 2 vl, vla, vlc, b, org \\
\end{xltabular}

\begin{sources}

\sourceitem%
  {A1}%
  {A-Ws}%
  {Cod. 735/11}%
  {autograph manuscript (principal source)}%
  {1822}%
  {}%
  {public domain}%
  {}%
  {full score; 14 pages}

\end{sources}


\textlt{\textit{Lyrics}\\Nos populus tuus et oves pascuae tuae,\\ confitebimur tibi Domine in saeculum,\\ in generationem et generationem\\ annuntiabimus laudem tuam.\\[1ex] (\bibleverse{Ps}(79/78:13))}

\section{Jubilate Deo · HerEy 77}

\begin{xltabular}{\linewidth}{@{} >\itshape l X}
genre & offertorium \\
festival & de Tempore \\
scoring & S, A, T, B (coro), 2 ob, 2 cl (C), 2 fag, 2 clno (C), 3 trb, timp (C–G), 2 vl, vla, vlc, b, org \\
\end{xltabular}

\begin{sources}

\sourceitem%
  {A1}%
  {A-Ws}%
  {Cod. 735/12}%
  {autograph manuscript (principal source)}%
  {1823}%
  {}%
  {public domain}%
  {}%
  {full score (20 pages); clno, trb, and timp written on 3 separate pages}

\end{sources}


\textlt{\textit{Lyrics}\\Jubilate Deo universa terra,\\ psalmum dicite nomini eius,\\ venite et audite et narrabo vobis,\\ omnes qui timetis Deum,\\ quanta fecit Dominus animae meae.\\ (\bibleverse{Ps}(66/65:1,2,16))}

\section{Confirma hoc Deus · HerEy 79}

\begin{xltabular}{\linewidth}{@{} >\itshape l X}
genre & offertorium \\
festival & in Festo Pentecostes \\
scoring & S, A, T, B (coro), 2 ob, 2 fag, 2 clno (C), timp (C–G), 2 vl, vla, vlc, b, org \\
\end{xltabular}

\begin{sources}

\sourceitem%
  {A1}%
  {A-Ws}%
  {Cod. 735/13}%
  {autograph manuscript (principal source)}%
  {1825}%
  {}%
  {public domain}%
  {}%
  {full score; 15 pages}

\end{sources}


\textlt{\textit{Lyrics}\\Confirma hoc Deus quod operatus es in nobis.\\ A templo sancto tuo quod est in Jerusalem tibi offerent reges munera.\\[1ex] (\bibleverse{Ps}(68/67:29-30))}

\section{Hæc est dies · HerEy 89}

\begin{xltabular}{\linewidth}{@{} >\itshape l X}
genre & offertorium \\
festival & in Festo Sanctæ Theresiæ \\
scoring & T (solo), S, A, T, B (coro), 2 ob, 2 clno (D), timp (D–A), 2 vl, vla, vla d’amore solo, vlc solo, b, org \\
\end{xltabular}

\begin{sources}

\sourceitem%
  {A1}%
  {A-Ws}%
  {Cod. 735/1}%
  {autograph manuscript (principal source)}%
  {1800-10}%
  {}%
  {public domain}%
  {}%
  {full score; 37 pages}

\end{sources}

\begin{xltabular}{\linewidth}{ll X}
\toprule
\itshape Bar & \itshape Staff & \itshape Description \\
\midrule \endhead
– & –   & Eybler reworked the two movements of this work into
          the offertory \textit{Summe Deus} (HerEy 90, 1818) and
          the gradual \textit{Te summe Jesu} (HerEy 49, 1809), respectively.
          To this end, he changed the lyrics and introduced several major
          changes (drafts of which also appear in \A1).
          Revision of first movement to HerEy~90:
          (1) Eybler deleted bars 61f and 91, replaced bars 96f with
          one new bar, and deleted bars~157f and 180–185,
          thereby shortening the movement from 212 to 200 bars.
          (2) He added parts for two bassoons and a mixed chorus.
          (3) He replaced the solo viola d’amore with a solo violin.
          Revision of second movement to HerEy~49:
          (1) Eybler replaced bars 232–234, 257f, and 278–291 with
          11, 8, and 36 bars, respectively, thereby extending
          the movement from 79 to 115 bars.
          (2) He added parts for two bassoons. \\
– & org & bass figures added by editor \\
1 & –   & tempo indication missing in \A1, here derived from HerEy 90 \\

\bottomrule
\end{xltabular}

\textlt{\textit{Lyrics}\\Haec est dies\\ qua candidae instar columbae\\ spiritus Theresiae\\ ad sacra templa caelitum se transtulit.\\[1ex] Sponsique voces audiit:\\ Veni soror de vertice carmeli,\\ ad agni nuptias veni,\\ ad gloriae coronam veni.\\[1ex] Te sponse Jesu,\\ te sponse virginum\\ beati ordines adorent\\ et nuptiali cantico laudent\\ per cuncta saecula.\\[1ex] (Cantus Index a01576)}

\section{Summe Deus · HerEy 90}

\begin{xltabular}{\linewidth}{@{} >\itshape l X}
genre & offertorium \\
festival & de Tempore \\
scoring & T (solo), S, A, T, B (coro), 2 ob, 2 fag, 2 clno (D), timp (D–A), 2 vl, vl solo, vla, vlc solo, b, org \\
\end{xltabular}

\begin{sources}

\sourceitem%
  {A1}%
  {A-Ws}%
  {Cod. 735/8}%
  {autograph manuscript (principal source)}%
  {1818}%
  {}%
  {public domain}%
  {}%
  {full score; 28 pages}

\end{sources}

\begin{xltabular}{\linewidth}{ll X}
\toprule
\itshape Bar & \itshape Staff & \itshape Description \\
\midrule \endhead
– & –   & This work is a revised version of the first movement
          of \textit{Hæc est dies} (HerEy 89). See the critical report
          of the latter work for details. \\
– & org & bass figures added by editor \\

\bottomrule
\end{xltabular}

\textlt{\textit{Lyrics}\\Summe Deus, te semper laudum carmine\\ voces nostrae celebrent\\ conjunctae angelorum elatis canticis.\\ Sint cordis nostri vota grata tibi,\\ sint oris nostri laudes carae tibi,\\ et pietatis tuae descendant in nos dona,\\ pie Pater, haec speramus.\\ Quam felix qui tanto\\ sub Patre coelestem tendit ad metam.}

\section{Magna et mirabilia · HerEy 108}

\begin{xltabular}{\linewidth}{@{} >\itshape l X}
genre & offertorium \\
festival & de Tempore \\
scoring & S, A, T, B (coro), 2 ob, 2 cl (\flat B), 2 fag, 2 cor (\flat B), 2 clno (\flat B), timp (\flat B–F), 2 vl, vla, vlc, b, org \\
\end{xltabular}

\begin{sources}

\sourceitem%
  {A1}%
  {A-Ws}%
  {Cod. 735/2}%
  {autograph manuscript (principal source)}%
  {1828}%
  {}%
  {public domain}%
  {}%
  {full score; 15 pages}


\sourceitem%
  {C1}%
  {A-Wn}%
  {HK.2153}%
  {print}%
  {1832}%
  {}%
  {public domain}%
  {https://data.onb.ac.at/rec/AC14328847}%
  {full score; Tobias Haslinger, Wien, plate number 5742}


\sourceitem%
  {C2}%
  {A-Wn}%
  {SA.87.C.106/1-2}%
  {print}%
  {1832}%
  {}%
  {public domain}%
  {https://data.onb.ac.at/rec/AC09173363}%
  {20 parts (S, A, T, B, ob 1, ob 2, cl 1, cl 2, fag 1, fag 2, cor 1, cor 2, clno 1, clno 2, timp, vl 1, vl 2, vla, vlc/b, org); Tobias Haslinger, Wien, plate number 5745}

\end{sources}

\begin{xltabular}{\linewidth}{ll X}
\toprule
\itshape Bar & \itshape Staff & \itshape Description \\
\midrule \endhead
105–110 & cor & ties only appear in \C1 and \C2 \\

\bottomrule
\end{xltabular}

\textlt{\textit{Lyrics}\\Magna et mirabilia sunt opera tua Domine.\\ Quis non timebit te,\\ et magnificabit nomen tuum?\\ Omnes gentes venient\\ et adorabunt in conspectu tuo.\\[1ex] (\bibleverse{Rev}(15:3,4))}


\section{Te Deum · HerEy 115}

\begin{xltabular}{\linewidth}{@{} >\itshape l X}
genre & hymn \\
festival & – \\
scoring & S, A, T, B (coro 1), S, A, T, B (coro 2), 2 fl, 2 ob, 2 cl (C), 2 fag, 2 cor (C), 4 clno (C), 3 trb, timp (C–G), 2 vl, vla, vlc, b, org \\
\end{xltabular}

\begin{sources}

\sourceitem%
  {A1}%
  {A-Ws}%
  {700}%
  {autograph manuscript (principal source)}%
  {1814-02}%
  {}%
  {public domain}%
  {}%
  {full score; 66 pages}

\end{sources}

\begin{xltabular}{\linewidth}{ll X}
\toprule
\itshape Bar & \itshape Staff & \itshape Description \\
\midrule \endhead
99 & B 2 & 1st \halfNote\ in \A1: e2 \\

\bottomrule
\end{xltabular}


\eesToc{}

\cleardoublepage%
\pagenumbering{arabic}%
\setcounter{page}{1}%
\includepdf[pages=-,link=true,linkname=score]{../../tmp/B4/full_score.pdf}%
\includepdf[pages=-,link=true,linkname=score,openright]{../../tmp/B4a/full_score.pdf}%

\end{document}
